
(i). D'autre part, (iii) implique (ii) gr\^ace \`a 2.3 et 2.8.

\subsection*{4. Th\'eor\`emes de finitude pour @@M0@@}

\textbf{4.1.} Soient @@M1@@ un sch\'ema et @@M2@@ un foncteur en groupes commutatifs sur @@M3@@. On d\'esigne par @@M4@@ (resp. @@M5@@) le sous-foncteur en groupes dont les @@M6@@-points @@M7@@ sont ceux de @@M8@@ tels que pour tout corps alg\'ebriquement clos @@M9@@ et tout @@M10@@-point @@M11@@ de @@M12@@, il existe un @@M13@@-sch\'ema connexe @@M14@@, deux @@M15@@-points @@M16@@ de @@M17@@ et un morphisme @@M18@@ tels que @@M19@@ et @@M20@@ dans @@M21@@ (resp. il existe un entier non-nul @@M22@@ tel que @@M23@@).

Il est \'evident que la formation de @@M24@@ (resp. @@M25@@) commute \`a tout changement de base, qu'un morphisme de foncteurs en groupes @@M26@@ induit un morphisme @@M27@@ (resp. @@M28@@), et que si @@M29@@ est un corps et @@M30@@ est repr\'esent\'e par un sch\'ema en groupes localement de type fini, alors @@M31@@ est repr\'esent\'e par la composante connexe de l'\'el\'ement neutre, qui est un sous-sch\'ema ouvert (et ferm\'e), de type fini, et g\'eom\'etriquement irr\'eductible (SGA 3 VI@@M32@@ 2.4).

\textbf{Lemme 4.2.} \emph{Soit @@M33@@ un morphisme entre deux foncteurs en groupes commutatifs sur un sch\'ema @@M34@@. On suppose que pour tout corps alg\'ebriquement clos @@M35@@ et tout morphisme @@M36@@, @@M37@@ (resp. @@M38@@) soit repr\'esentable par un sch\'ema en groupes localement de type fini, et @@M39@@ soit repr\'esentable par un morphisme de type fini, (c-\`a-d, quasi-compact). Alors}
@@M40@@

